\subsection{域上的有限型代数}

本节中，k 表示一个域。

引理 3. --- 设 A 是一个有限生成 k-代数，\(\mathfrak{p}_{0} \subset \ldots \subset \mathfrak{p}_{m}\) 是 A 的一条素理想极大链。存在一个整数 \(n \geqslant m\)、A 的元素序列 \((x_{1}, \ldots, x_{n})\)，它们在 \(k\) 上代数自由（A, IV, p. 4），并且满足：

\begin{enumerate}
\def\labelenumi{\alph{enumi})}
\item
  A 在环 \(\mathbf{B} = k[x_1, \ldots, x_n]\) 上整；
\item
  对每个满足 \(0 \leqslant j \leqslant m\) 的 j，理想 \(\mathfrak{p}_{j} \cap B\) 由 \(x_{k}\)（\(1 \leqslant k \leqslant n - m + j\)）生成。
\end{enumerate}

根据归一化引理（V, § 3, no. 1, th. 1），存在一个整数 \(n \geqslant 0\)、A 的代数自由元素的有限序列 \((x_1, ..., x_n)\)，以及一个递增整数序列 \((h(j))_{0 \leqslant j \leqslant m}\)，其各项 \(\leqslant n\)，使得理想 \(\mathfrak{p}_j \cap \mathbf{B}\) 等于 B 的素理想 \(\mathfrak{q}_j\)，后者由 \(x_k\)（\(1 \leqslant k \leqslant h(j)\)）生成，并且 A 在环 B 上整。设 \(j\) 是满足 \(0 \leqslant j < m\) 的整数。通过取商，由 B 到 A 的典范嵌入得到从 \(\mathrm{B} / \mathfrak{q}_j\) 到 \(\mathrm{A} / \mathfrak{p}_j\) 的单射，它使 \(\mathrm{A} / \mathfrak{p}_j\) 成为 \((\mathrm{B} / \mathfrak{q}_j)\)-有限代数。由于环 \(\mathrm{B} / \mathfrak{q}_j\) 同构于含有 \(n - h(j)\) 个不定元的、在 \(k\) 上的多项式代数，故它是整闭的（V, § 1, no. 3, 命题 13 的推论 3）。根据 no. 3 的定理 2，我们因此有

\[
1 = \operatorname{ht} \left(\mathfrak {p} _ {j + 1} / \mathfrak {p} _ {j}\right) = \operatorname{ht} \left(\mathfrak {q} _ {j + 1} / \mathfrak {q} _ {j}\right) \geqslant h (j + 1) - h (j).
\]

由此有 \(h(j + 1) \leqslant h(j) + 1\) 且 \(\mathfrak{q}_{j+1} \neq \mathfrak{q}_j\)，从而 \(h(j + 1) = h(j) + 1\)。但是，有 \(h(m) = n\)，因为 \(\mathfrak{q}_m\) 是极大理想（V, § 2, no. 1, 命题 1），最后得到 \(h(j) = n - m + j\)。

定理 3. — 设 A 是一个有限生成 k-代数。

\begin{enumerate}
\def\labelenumi{\alph{enumi})}
\item
  对 A 的每个极小素理想  ，从  开始的 A 的素理想极大链的长度都等于整数 \_。
\item
  环 A 是串联的，其维数是整数 \(\deg. \operatorname{tr}_{k}\kappa(\mathfrak{p})\) 的上确界，其中 \(\mathfrak{p}\) 遍历 A 的极小素理想。
\item
  如果 A 是一个整环，则 A 的所有素理想极大链具有相同长度，并且 A 的维数是 A 的分式域关于 k 的超越次数。
\end{enumerate}

假设 A 是整的，考虑 A 的一条素理想极大链 \(\mathfrak{p}_0 \subset \ldots \subset \mathfrak{p}_m\)。有 \(\mathfrak{p}_0 = 0\)。于是由引理 3 得到一个单射 \(\varphi : k[\mathrm{X}_1, \ldots, \mathrm{X}_m] \to \mathrm{A}\)，它是 \(k\)-代数同态，并使 A 成为 \(k[\mathrm{X}_1, \ldots, \mathrm{X}_m]\) 上的有限代数。因此，A 的分式域关于 \(k\) 的超越次数等于 \(m\)，从而 \(c\)。断言 \(a)\) 由将 \(c)\) 应用于环  得出，而断言 \(b)\) 立即由 \(a)\) 及 § 1, no. 2 的命题 5 得出。

推论 1. --- 设 n 是一个正整数。则

\[
\dim (k [ \mathrm{X} _ {1},..., \mathrm{X} _ {n} ]) = n.
\]

对于有限型 \(k\)-代数 A，要使其维数为 \(n\)，必要且充分条件是存在一个 \(k\)-同态 \(\varphi: k[X_1, ..., X_n] \to A\)，使 A 成为 \(k[X_1, ..., X_n]\) 上的有限代数。

这由定理 3、归一化引理（V, § 3, no. 1, 定理 1）以及 no. 3 的定理 1 a) 得出。

推论 2. --- 设 A 是 k 上的有限生成整环。对于 A 的每个素理想  ，有

\[
\begin{array}{r l} \mathrm{ht} (\mathfrak {p}) & = \dim (\mathrm{A} _ {\mathfrak {p}}) = \dim (\mathrm{A}) - \dim (\mathrm{A} / \mathfrak {p}) \\ & = \dim (\mathrm{A}) - \deg . \operatorname{tr} _ {k} \kappa (\mathfrak {p}). \end{array}
\]

特别地，对 A 的每个极大理想 \_，都有  。这由定理 3 及 § 1, no. 3 的注 4 得出。

推论 3. --- 设 A 是一个有限生成 k-代数，且 f 是不属于 A 的任何极小素理想的 A 元素（例如 f 不是零因子，参见 IV, § 1, no. 1, 命题 2 的推论 3 及 no. 3 的推论 1）。则有 \(_{f}\)。

映射 \(\mathfrak{p} \mapsto \mathfrak{p}\mathrm{A}_f\) 是从 A 的极小素理想集合到 \(\mathrm{A}_f\) 的极小素理想集合的双射。此外，环 \(\mathrm{A} / \mathfrak{p}\) 与 \(\mathrm{A}_f / \mathfrak{p}\mathrm{A}_f = (\mathrm{A} / \mathfrak{p})_f\) 具有相同的分式域。因此，只需应用定理 3 b)。

推论 4. --- 设 A 是一个有限生成 k-代数，且  是 A 的一个素理想。a) 要使  成为极大理想，必要且充分条件是  的分式域是 k 的有限扩域。

\begin{enumerate}
\def\labelenumi{\alph{enumi})}
\setcounter{enumi}{1}
\tightlist
\item
  设 \(f \in \mathbf{A} - \mathfrak{p}\)。当且仅当 \(\mathfrak{p}\) 是 \(\mathbf{A}\) 的极大理想时，\(\mathfrak{p}\mathbf{A}_f\) 是 \(\mathbf{A}_f\) 的极大理想。
\end{enumerate}

如果  是 A 的极大理想，则 A/p 是域，因而是维数 0 的环；它是 k 的有限生成扩域且其超越次数为 0（定理 3 c)），所以它是 k 的有限扩域。反之，如果 A/p 的分式域是 k 的有限扩域，则 \_，所以  是极大理想。断言 b) 由断言 a) 以及 A/p 与 \_ 具有相同分式域这一事实得出。

推论 4 的断言 \(a\)）是零点定理的一种形式（V, § 3, no. 3, 命题 1）。

推论 5. --- 设 A 是一个有限生成 k-代数，  是 A 的一个素理想，\((\mathfrak{p}_{i})_{i\in I}\) 是包含于  中的 A 的极小素理想族。我们有：

\[
\begin{array}{r l} \dim_ {\mathfrak {p}} (\mathrm{A}) & = \sup _ {i \in \mathrm{I}} \dim (\mathrm{A} / \mathfrak {p} _ {i}) \\ & = \dim (\mathrm{A} _ {\mathfrak {p}}) + \dim (\mathrm{A} / \mathfrak {p}) \\ & = \dim (\mathrm{A} _ {\mathfrak {p}}) + \deg . \operatorname{tr} _ {k} \kappa (\mathfrak {p}). \end{array}
\]

我们有 \(\dim_{\mathfrak{p}}(\mathbf{A}) = \sup_{i\in I}\dim_{\mathfrak{p} / \mathfrak{p}_i}(\mathbf{A} / \mathfrak{p}_i)\)。根据 § 1, no. 3 的命题 6。但是，根据推论 3，有 \(\dim_{\mathfrak{p} / \mathfrak{p}_i}(\mathbf{A} / \mathfrak{p}_i) = \dim (\mathbf{A} / \mathfrak{p}_i)\)，从而得到第一个等式。根据推论 2，有 \(\dim (\mathbf{A} / \mathfrak{p}_i) = \dim ((\mathbf{A} / \mathfrak{p}_i)_{\mathfrak{p}}) + \dim (\mathbf{A} / \mathfrak{p})\)。因此，推论的第二个等式由 \(\dim (\mathbf{A}_{\mathfrak{p}}) = \sup_{i\in I}\dim ((\mathbf{A} / \mathfrak{p}_i)_{\mathfrak{p}})\) 这一事实以及定理 3 的第三个结论得出。

因此，\(\dim_{\mathfrak{p}}(\mathbf{A})\) 是包含于 \(\mathfrak{p}\) 中的 A 的素理想链的长度的上确界。

推论 6. --- 设 A 是一个不等于 0 的有限生成 k-代数，n 是一个整数 。下列条件等价：

\begin{enumerate}
\def\labelenumi{\alph{enumi})}
\item
  对每个 \(\mathfrak{p} \in \operatorname{Ass}(\mathbf{A})\)，有 \(\dim(\mathbf{A}/\mathfrak{p}) = n\)。
\item
  与 A 相关的每个素理想都是极小的，并且 Spec(A) 的所有不可约分支的维数都是 n。
\item
存在一个单射 \(k\)-同态 \(\varphi: k[X_1, ..., X_n] \to A\)，使 \(A\) 成为 \(k[X_1, ..., X_n]\) 上的有限型且无挠的模。
\end{enumerate}

\(a)\) 的等价性是显然的。证明 \(b)\) 蕴含 \(a)\)。由 \(c)\)\(b)\)，环 A 的维数为 \(n\)，因而存在（推论 1）一个 \(k\)-同态 \(\varphi : k[X_1, ..., X_n] \to A\)，使 A 成为 \(k[X_1, ..., X_n]\) 上的有限型模。对于每个 \(\mathfrak{p} \in \operatorname{Ass}(A)\)，环 \(\mathfrak{p}\) 在 \(k[X_1, ..., X_n]\) 上是整的，因而有 \(n = \dim(A/\mathfrak{p}) = \dim(k[X_1, ..., X_n]/\varphi^{-1}(\mathfrak{p}))\)，根据定理 1 的 \(a)\)（no. 3），从而 \(\varphi^{-1}(\mathfrak{p}) = 0\)。单射同态 \(\varphi\) 的 \(k[X_1, ..., X_n]\) 非零元素的像不是 A 中的零因子（IV, § 1, no. 1, 命题 2 的推论 3），因此得到 \(c)\)。

反之，假设条件 \(c\) 成立。对于每个 \(\mathfrak{p} \in \operatorname{Ass}(A)\)，由 \(k[X_1, ..., X_n] \to A/\mathfrak{p}\) 诱导的同态 \(\varphi\) 是单射（同上）。因此，根据推论 1，有 \(\dim(A/\mathfrak{p}) = n\)。

由 \(a)\)、\(b)\)、\(c)\) 可知，\(\dim_{\mathfrak{p}}(\mathbf{A}) = \dim(\mathbf{A})\) 对每个 \(\mathfrak{p}\) 和 \(\mathbf{A}\) 成立。

命题 4. --- 设 A、B 是两个有限型 k-代数，\(\rho: \mathbf{A} \to \mathbf{B}\) 是一个代数同态。假设 A 是一个整环，B 作为 A-模无挠，并令 K 为 A 的分式域。我们有

\[
\dim (\mathbf {B}) = \dim (\mathbf {A}) + \dim (\mathbf {B} \otimes_ {\mathbf {A}} \mathbf {K}).
\]

先假设 B 是一个整环。代数 B \(\otimes_{\mathbf{A}} \mathbf{K}\) 是由一个不含 0 的乘法子集定义的 B 的一个分式环。因此，它的分式域是 B 的分式域 L。根据定理 3，有

\[
\dim (\mathbf {B}) = \deg . \operatorname{tr} _ {k} \mathrm{L}, \quad \dim (\mathrm{A}) = \deg . \operatorname{tr} _ {k} \mathrm{K},
\]

\[
\dim (\mathbf {B} \otimes_ {\mathbf {A}} \mathbf {K}) = \deg . \operatorname{tr} _ {\mathbf {K}} \mathbf {L}.
\]

现在根据 A, V, p. 106 中的推论，有

\[
\deg . \operatorname{tr} _ {k} \mathrm{L} = \deg . \operatorname{tr} _ {\mathrm{K}} \mathrm{L} + \deg . \operatorname{tr} _ {k} \mathrm{K},
\]

从而在此情形下得到结论。

下面转到一般情形。B 的每个极小素理想  都由 B 中的零因子组成，因而位于 A 的理想 0 之上。因此，映射 \(\mathfrak{p} \mapsto \mathfrak{p}.(\mathbf{B} \otimes_{\mathbf{A}} \mathbf{K})\) 是从 B 的极小素理想集合到 \(B \otimes_{A} K\) 的极小素理想集合的双射。命题由证明的第一部分以及 § 1, no. 3 的命题 6 c) 得出。

推论。--- 设 \(\rho: A \to B\) 是有限型 \(k\)-代数的一个单射同态。则有 \(\dim(A) \leqslant \dim(B)\)。

事实上，设  是 A 的一个极小素理想，使得  。存在 B 的一个位于  之上的素理想  （II, § 2, no. 6, 命题 16）。将命题 4 应用于 \(\dim(A) = \dim(A/p)\) 和  ，有 \(\dim(A) = \dim(A/\mathfrak{p}) \leqslant \dim(B/\mathfrak{q}) \leqslant \dim(B)\)，从而得到推论。

引理 4. --- 设 A 和 B 是 k-代数的两个整环，M 是无挠 A-模，N 是无挠 B-模。如果环  是整环，则 \(_{k}\) 是关于 \(_{k}\)\(_{k}\) 的无挠模。

设 K（分别为 L）是 A（分别为 B）的分式域。存在一个集合 I（分别为 J），使得 M（分别为 N）同构于 \(\mathbf{K}^{(\mathbf{I})}\)（分别为 \(\mathbf{L}^{(\mathbf{J})}\)）的一个子模。于是， -模 \((\mathbf{A} \otimes_k \mathbf{B})\) 同构于 \(M \otimes_k N\) 的一个子模，而后者同构于 \(\mathbf{K}^{(\mathbf{I})} \otimes_k \mathbf{L}^{(\mathbf{J})}\)。由于 \((\mathbf{K} \otimes_k \mathbf{L})^{(\mathbf{I} \times \mathbf{J})}\) 是整环 \(K \otimes_k L\) 的一个分式环，所以它作为 \(A \otimes_k B\)\(A \otimes_k B\) 上的模是无挠的，从而得到引理。

命题 5. --- 设 \(k'\) 是 \(k\) 的扩域，A 是有限型 \(k\)-代数，B 是有限型 \(k'\)-代数。

\begin{enumerate}
\def\labelenumi{\alph{enumi})}
\tightlist
\item
  \(k'\)-代数 \(\mathbf{A} \otimes_k \mathbf{B}\) 是有限型的，并且有
\end{enumerate}

\[
\dim (\mathbf {A} \otimes_ {k} \mathbf {B}) = \dim (\mathbf {A}) + \dim (\mathbf {B}).
\]

\begin{enumerate}
\def\labelenumi{\alph{enumi})}
\setcounter{enumi}{1}
\tightlist
\item
  设  是  的一个素理想；设  （分别为  ）是 \(\otimes_{k}\) 在 A（分别为 B）中的逆像。我们有
\end{enumerate}

\[
\dim_ {\mathfrak {r}} (\mathrm{A} \otimes_ {k} \mathrm{B}) = \dim_ {\mathfrak {p}} (\mathrm{A}) + \dim_ {\mathfrak {q}} (\mathrm{B}).
\]

置 \(n = \dim(A)\)，\(m = \dim(B)\)。根据定理 3 的推论 1，存在单射代数同态 \(\varphi: k[X_1, ..., X_n] \to A\) 和 \(\psi: k'[Y_1, ..., Y_m] \to B\)，使 \(A\) 和 \(B\) 分别成为 \(k[X_1, ..., X_n]\) 和 \(k'[Y_1, ..., Y_m]\) 上的有限代数。于是同态 \(\varphi \otimes \psi\) 是单射，并使 \(A \otimes_k B\) 成为一个 \(k'\)-代数 \(k[X_1, ..., X_n] \otimes_k k'[Y_1, ..., Y_m]\) 上的有限代数，后者可识别为 \(k'[X_1, ..., X_n, Y_1, ..., Y_m]\)。因此，根据定理 3 的推论 1，有 \(\dim(A \otimes_k B) = n + m\)，这证明了 \(a\)。

注意，当 A 和 B 是整环时，\(A \otimes_k B\) 是无挠的非零有限生成 \(k'[\mathrm{X}_{1}, \ldots, \mathrm{X}_{n}, \mathrm{Y}_{1}, \ldots, \mathrm{Y}_{m}]\)-模（根据引理 4），因而有

\[
\dim_ {\mathbf {r}} (\mathrm{A} \otimes_ {k} \mathrm{B}) = n + m = \dim (\mathrm{A}) + \dim (\mathrm{B})
\]

  根据注 1，对于 \(\otimes_{k}\) 的每个素理想  ，均有

现在证明 \(b\)。设 \(\mathfrak{r}_{0}\) 是 \(\otimes_{k}\) 的一个包含于 \(\mathfrak{r}\) 的极小素理想，并记 \(\mathfrak{p}_{0}\)（分别为 \(\mathfrak{q}_{0}\)）为 \(\mathfrak{r}_{0}\) 在 A（分别在 B）中的逆像。环  同构于环 \((\mathbf{A} \otimes_k \mathbf{B}) / \mathfrak{r}_{0}\) 的一个商环 \((\mathbf{A} / \mathfrak{p}_{0}) \otimes_k (\mathbf{B} / \mathfrak{q}_{0})\)。因此，根据 \(a\)，有

\[
\dim \left(\left(\mathrm{A} \otimes_ {k} \mathrm{B}\right) / \mathfrak {r} _ {0}\right) \leqslant \dim (\mathrm{A} / \mathfrak {p} _ {0}) + \dim (\mathrm{B} / \mathfrak {q} _ {0}).
\]

应用定理 3 的推论 5，我们得到不等式

\[
\dim_ {\mathfrak {r}} (\mathrm{A} \otimes_ {k} \mathrm{B}) \leqslant \dim_ {\mathfrak {p}} (\mathrm{A}) + \dim_ {\mathfrak {q}} (\mathrm{B}).
\]

反之，设  （分别为 ）是 A（分别为 B）中包含于 \(\mathfrak{p}_{0}\)（分别为 \(q_{0}\)）的一个极小素理想。根据上面的注记，有

\[
\dim (\mathrm{A} / \mathfrak {p} _ {0}) + \dim (\mathrm{B} / \mathfrak {q} _ {0}) = \dim_ {\overline {{\mathfrak {r}}}} ((\mathrm{A} / \mathfrak {p} _ {0}) \otimes_ {k} (\mathrm{B} / \mathfrak {q} _ {0}))
\]

其中 \(\overline{\mathbf{r}}\) 是  在典范满射 \(\mathbf{A} \otimes_{k} \mathbf{B} \to (\mathbf{A}/\mathfrak{p}_{0}) \otimes_{k} (\mathbf{B}/\mathfrak{q}_{0})\) 下的像。前一等式的第二项显然小于 \(\dim_{\mathfrak{r}}(\mathbf{A} \otimes_{k} \mathbf{B})\)。应用定理 3 的推论 5，我们得到不等式

\[
\dim_ {\mathfrak {p}} (\mathrm{A}) + \dim_ {\mathfrak {q}} (\mathrm{B}) \leqslant \dim_ {\mathfrak {r}} (\mathrm{A} \otimes_ {k} \mathrm{B}),
\]

这就完成了证明。

推论。--- 设 A 是有限型 \(k\)-代数，\(k'\) 是 \(k\) 的扩域，且 \(\mathbf{A}'\) 是 \(k'\)-代数 \(\mathbf{A} \otimes_k k'\)。

\begin{enumerate}
\def\labelenumi{\alph{enumi})}
\item
  有 \(\dim (\mathbf{A}^{\prime}) = \dim (\mathbf{A})\)。
\item
  设 \(\mathfrak{p}'\) 是 \(\mathbf{A}'\) 的一个素理想，\(\mathfrak{p}\) 是它在 \(\mathbf{A}\) 中的逆像；有 \(\dim_{\mathfrak{p}'}(\mathbf{A}') = \dim_{\mathfrak{p}}(\mathbf{A})\)。
\item
  设 \(\mathfrak{p}'\) 是 \(\mathbf{A}'\) 的一个极小素理想，\(\mathfrak{p}\) 是它在 \(\mathbf{A}\) 中的逆像。则 \(\mathfrak{p}\) 是极小的，并且有 \(\dim(\mathbf{A}'/\mathfrak{p}') = \dim(\mathbf{A}/\mathfrak{p})\)
\end{enumerate}

由命题 5 可得 \(a)\) 和 \(b)\)（其中取 \(\mathbf{B} = k'\)）。下面证明 \(c)\)。理想  是极小的（no. 1, 命题 1），并且有

\[
\dim (\mathrm{A} ^ {\prime} / \mathfrak {p} ^ {\prime}) = \dim_ {\mathfrak {p} ^ {\prime}} (\mathrm{A} ^ {\prime}) = \dim_ {\mathfrak {p}} (\mathrm{A}) = \dim (\mathrm{A} / \mathfrak {p}).
\]

注 2. --- 假设  是  的一个纯不可分扩张。则典范映射  是同胚。

事实上，设 \(\mathfrak{p} \in \operatorname{Spec}(\mathbf{A})\)。根据 § 1 的引理 1，空间 \(f^{-1}(\{\mathfrak{p}\})\) 同胚于 \(\operatorname{Spec}(\kappa(\mathfrak{p}) \otimes_k k')\)。现在，集合 \(\alpha\) 是 \(\kappa(\mathfrak{p}) \otimes_k k'\) 的幂零元集合，是一个素理想（A, V, p. 134, 推论），而商环 \(\kappa(\mathfrak{p}) \otimes_k k')/\mathfrak{a}\) 是一个在 \(\kappa(\mathfrak{p})\) 上整的整环，因而是一个域（A, V, p. 16, 推论 1 及 p. 10, 命题 1）。因此 \(f^{-1}(\{\mathfrak{p}\})\) 退化为单个元素。由此可知，映射 \(f\) 是从 \(\operatorname{Spec}(\mathbf{A}')\) 到 \(\operatorname{Spec}(\mathbf{A})\) 的递增双射；这两个集合按素理想的包含关系排序，因而在 \(\operatorname{Spec}(\mathbf{A})\) 与 \(\operatorname{Spec}(\mathbf{A}')\) 的不可约闭子集之间诱导出双射。由于 \(\operatorname{Spec}(\mathbf{A})\)（分别为 \(\operatorname{Spec}(\mathbf{A}')\)）的闭子集都是不可约闭子集的有限并，\(f\) 是同胚。

关于一般情形，参见第 98 页练习 24。

